\chapter{Background and Related Work}

% =====================================================================
% Dos mitades: fundamentos necesarios para seguir la propuesta (secciones 1-3)
% y estado del arte que la justifica (secciones 4-6). La sección final es la
% bisagra hacia el Capítulo 3 y no debe omitirse.
% =====================================================================

\section{Vehicular Telemetry: Data Sources and Service Requirements}
\emph{Caracterizar el tipo de datos que genera un vehículo instrumentado y sus
requisitos diferenciados de fiabilidad, periodicidad y vigencia temporal.
Establecer aquí la noción de plazo de validez del dato, que se usará en toda la
tesis.}

\section{Heterogeneous Vehicular Access: Cellular 5G NR and IEEE 802.11}
\emph{Comparar ambas tecnologías de acceso en los términos relevantes para el
problema: cobertura, disponibilidad, capacidad y coste energético. El objetivo no
es describirlas exhaustivamente, sino establecer que son complementarias y que
esa complementariedad es explotable.}

\section{Energy Consumption of Wireless Interfaces in Vehicular Nodes}
\emph{Revisar cómo se modela el consumo energético de interfaces inalámbricas en
la literatura y en las herramientas de simulación existentes, distinguiendo
modelos basados en estado del transceptor frente a modelos basados en actividad.
Esta sección motiva el modelo energético propio del Capítulo 4.}

\section{Opportunistic Offloading and Connectivity Prediction}
\emph{Estado del arte del núcleo de la tesis: esquemas de descarga oportunista,
retención de datos a la espera de conectividad y predicción de contactos.
Clasificar los trabajos según qué información usan para decidir y qué garantizan.}

\section{Deadline-Aware and Freshness-Constrained Data Delivery}
\emph{Revisar los trabajos que incorporan plazos o métricas de frescura de la
información como restricción de la decisión de transmisión, y señalar en qué
medida se combinan con criterios energéticos.}

\section{Simulation Environments for Heterogeneous Vehicular Networks}
\emph{Revisar las plataformas disponibles para evaluar escenarios vehiculares
heterogéneos y su grado de soporte a la doble interfaz y al modelado energético.
Justificar la combinación adoptada sin entrar en detalles de integración, que
corresponden al Capítulo 4.}

\section{Research Gap and Positioning of This Work}
\emph{Sección bisagra. Sintetizar qué queda sin resolver en las secciones
anteriores —típicamente: predicción de oportunidades, restricción de memoria
finita, plazo de validez y coste energético tratados conjuntamente y validados
sobre movilidad realista— y enunciar cómo esta tesis ocupa ese hueco. Debe
enlazar directamente con el capítulo siguiente.}
